\setchapterpreamble[u]{\margintoc}
\chapter{T-SQL: Procedimentos Armazenados e Transações}
\labch{cet615-procedimentos}

Este capítulo cobre a última apresentação da segunda
aula\sidecite{BaiocoCET615}: procedimentos armazenados e o uso deles
para encapsular transações. É o conteúdo da questão 2 da primeira
avaliação (\refch{cet615-avaliacoes}).

\section{O que é um procedimento armazenado}

Um \textbf{procedimento armazenado}\index{procedimento armazenado}
(\emph{stored procedure}) é um conjunto de comandos SQL compilado e
guardado no servidor, dentro de um banco. A analogia da aula é a de uma
rotina de linguagem de programação: a lógica fica no banco, e a
aplicação só chama procedimentos e visões, com o mínimo de SQL escrito
nela.

Vantagens:

\begin{itemize}
	\item \textbf{Encapsulamento.} Rotinas frequentes ficam disponíveis
	para todas as aplicações, em vez de reescritas em cada uma.
	\item \textbf{Menos erros.} Em tese, os procedimentos já foram
	testados. Em aula, a professora disse esperar que sim na prática, e
	que pelo menos os erros ficam mais fáceis de localizar.
	\item \textbf{Segurança.} O acesso fica limitado ao que o
	procedimento faz, não importa quantas tabelas ele use. Com visões e
	procedimentos, a permissão é dada por funcionalidade, como a interface
	libera ou bloqueia uma opção, e não por tabela.
\end{itemize}

No SSMS em português, os procedimentos ficam em \emph{Programação}
$\rightarrow$ \emph{Procedimentos Armazenados}.

\section{Criar e executar}

\begin{lstlisting}[caption={Busca de livros por parte do título.}]
create procedure busca_livro
@trecho varchar(60)
as
select codlivro, titulo
from livro
where titulo like '%' + @trecho + '%'
go

exec busca_livro 'casmurro'
execute busca_livro @trecho = 'casmurro'   -- passagem por nome
\end{lstlisting}

O parâmetro é \texttt{varchar} de propósito. Um \texttt{char(60)} viria
completado com espaços à direita, e \texttt{'\%' + @trecho + '\%'}
procuraria o trecho seguido de brancos. Os parâmetros que vão para
colunas, porém, devem ter o mesmo tipo da coluna.

A passagem pode ser posicional ou por nome. Depois que a chamada começa
a passar parâmetros por nome, os seguintes também precisam ser por nome.
Um \texttt{varchar} vai entre aspas na chamada.

Quando a consulta tem vários parâmetros (por exemplo, montar o padrão do
\texttt{like}), a professora considera o procedimento melhor que a
visão: ele recebe os valores e monta o \texttt{where} por dentro, em vez
de exigir um \texttt{where} de quem consulta a visão.

\subsection{Valor padrão}

\begin{lstlisting}
create procedure busca_livro2
@trecho varchar(60) = '%'
as
select codlivro, titulo from livro where titulo like @trecho
go

exec busca_livro2            -- sem argumento: traz todos
exec busca_livro2 'Dom%'
\end{lstlisting}

É por causa de valores padrão que alguns procedimentos do sistema
funcionam com ou sem parâmetro.

\subsection{Parâmetro de saída}

\begin{lstlisting}
create procedure titulo_do_livro
@codlivro int,
@titulo char(60) output       -- devolvido a quem chamou
as
select @titulo = titulo from livro where codlivro = @codlivro
go

declare @t char(60)
exec titulo_do_livro 10, @t output   -- output também na chamada
print @t
\end{lstlisting}

O parâmetro \texttt{output} é passado por referência. A ideia é não
mostrar o resultado na tela de dentro do procedimento e sim devolvê-lo
numa variável. A variável é atribuída no próprio \texttt{select}
(\texttt{select @titulo = titulo ...}). Quem chama declara uma variável
do mesmo tipo e escreve \texttt{output} também no \texttt{exec}. Sem o
\texttt{print}, nada aparece.

\subsection{Valor de retorno}

\begin{lstlisting}
create procedure conta_livros
@trecho varchar(60)
as
select codlivro, titulo from livro where titulo like '%' + @trecho + '%'
return @@rowcount
go

declare @r int
exec @r = conta_livros 'amor'
print @r
\end{lstlisting}

\begin{itemize}
	\item O \texttt{return} só devolve \textbf{inteiro}. Serve para status
	(sucesso ou falha), não para devolver o valor de uma coluna; para isso
	existe o \texttt{output}.
	\item O \texttt{return} encerra o procedimento ali mesmo, inclusive
	dentro de um \texttt{if}.
	\item O retorno é opcional: o procedimento funciona igual sem ele. Para
	usá-lo, declara-se uma variável e executa-se o bloco inteiro
	(\texttt{declare}, \texttt{exec @r = ...}, \texttt{print}).
	\item \texttt{print} aparece na aba \emph{Mensagens} do SSMS;
	\texttt{select @r} aparece na aba \emph{Resultados}, como uma
	consulta.
\end{itemize}

\subsection{Alterar e excluir}

\begin{lstlisting}
alter procedure busca_livro
@trecho varchar(60) = '%'
as
select codlivro, titulo from livro where titulo like '%' + @trecho + '%'
go

exec sp_helptext busca_livro           -- mostra o código
exec sp_rename 'busca_livro', 'procura_livro'
drop procedure procura_livro, busca_livro2
\end{lstlisting}

Como na visão, o \texttt{alter} substitui a definição e mantém as
permissões.

\section{Procedimentos com transação}

O uso principal na disciplina: encapsular uma transação num
procedimento que devolve 1 em caso de sucesso e 0 em caso de falha. O
padrão é o do \refch{cet615-dml}, com um \texttt{return} em cada ramo.

\begin{lstlisting}[caption={Cadastro de livro com exemplar: obra genérica e específica numa transação.}]
create procedure ins_livro_exemplar
@codlivro int,
@titulo char(60),
@isbn numeric(13,0),
@codautor int,
@codexemplar int
as
begin transaction
insert into livro (codlivro, titulo, isbn, codautor)
values (@codlivro, @titulo, @isbn, @codautor)
if @@rowcount > 0               /* livro inserido */
begin
   insert into exemplar (codexemplar, codlivro, estado)
   values (@codexemplar, @codlivro, 'N')   -- estado fixo: novo
   if @@rowcount > 0            /* exemplar inserido */
   begin
      commit transaction
      return 1
   end
   else
   begin
      rollback transaction
      return 0
   end
end
else
begin
   rollback transaction
   return 0
end
go

declare @ret int
exec @ret = ins_livro_exemplar 10, 'Dom Casmurro', 9788535910663, 1, 100
print @ret
\end{lstlisting}

\subsection{Como projetar os parâmetros}

A professora deu atenção especial a esse ponto:

\begin{itemize}
	\item \textbf{A chave vai uma vez só.} \texttt{@codlivro} serve para as
	duas tabelas. Passar o código duas vezes abre a chance de chegarem
	valores diferentes, um erro que ela provocou de propósito na transação
	sem procedimento.
	\item \textbf{O que é fixo não é parâmetro.} O tipo da especialização
	(0 para pessoa física, 0 para livro) e o estado inicial são escritos
	no próprio procedimento. Como parâmetro, só abririam espaço para erro.
	\item \textbf{Um procedimento por especialização.} Juntar o cadastro de
	pessoa física e jurídica num só (``a mania de querer aproveitar'')
	deixaria metade dos parâmetros nula em cada chamada. Para exclusão,
	um procedimento único faz sentido: basta o código, e o procedimento
	descobre o tipo.
	\item \textbf{A transação é a menor possível.} O que não precisa estar
	junto fica de fora. Por isso a avaliação diz que o vendedor já está
	cadastrado: cadastrar cliente é outra transação.
	\item \textbf{A ordem dos \texttt{insert} segue as chaves
	estrangeiras.} Primeiro a linha referenciada (a obra genérica), depois
	a que referencia (a específica).
\end{itemize}

\subsection{Por que um \texttt{return} em cada ramo}

Não dá para escrever um único \texttt{return @@rowcount} no fim, porque
o \texttt{@@rowcount} muda a cada comando. A documentação\sidecite{MSRowcount}
detalha: comandos DML gravam nele o número de linhas afetadas; comandos
como \texttt{print}, \texttt{begin transaction} e \texttt{commit
transaction} o zeram; atribuições simples o deixam em 1. Depois do
\texttt{commit}, ele já vale 0. A regra prática é testar o
\texttt{@@rowcount} logo depois do comando que interessa, ou guardá-lo
numa variável. Com \texttt{return} junto do \texttt{commit} ou do
\texttt{rollback}, o bloco \texttt{begin ... end} passa a ser
obrigatório em cada ramo.

\subsection{Procedimentos que chamam procedimentos}

Um procedimento pode executar outro. No exemplo da aula, um
procedimento de transação completa cadastra um cliente pessoa física e
duas contas (corrente e poupança) chamando os procedimentos menores e
testando o retorno de cada um:

\begin{lstlisting}
declare @retorno int
begin transaction
exec @retorno = ins_pessoa_fisica ...
if @retorno > 0
begin
   exec @retorno = ins_conta ... 0 ...   -- corrente
   if @retorno > 0
   begin
      exec @retorno = ins_conta ... 1 ...  -- poupança
      if @retorno > 0 commit transaction
      else rollback transaction
   end
   else rollback transaction
end
else rollback transaction
return @retorno
\end{lstlisting}

Aqui o \texttt{return @retorno} pode ficar no fim, porque
\texttt{@retorno} é uma variável comum e não muda sozinha. Modularizar
encurta o procedimento maior, e os menores já foram testados. Em
compensação, uma transação grande trava mais o banco, e a recomendação
continua sendo mantê-las enxutas.

\subsection*{Complemento: \texttt{rollback} dentro de procedimento aninhado}

A aula não tratou deste ponto, e ele afeta o exemplo acima no caminho de
falha. O SQL Server não tem transações aninhadas
independentes\sidecite{MSRollback}:

\begin{itemize}
	\item cada \texttt{begin transaction} soma 1 em
	\texttt{@@trancount}, e o \texttt{commit} de uma transação interna só
	subtrai 1. Quem efetiva de fato é o \texttt{commit} mais externo;
	\item um \texttt{rollback} sem ponto de salvamento desfaz tudo até o
	\texttt{begin transaction} mais externo e zera
	\texttt{@@trancount}, mesmo que tenha sido executado dentro do
	procedimento chamado;
	\item se o procedimento termina com \texttt{@@trancount} diferente do
	valor com que começou, o SQL Server emite uma mensagem informativa.
\end{itemize}

Na prática: se \texttt{ins\_conta} falha e executa \texttt{rollback}, a
transação inteira, inclusive o cliente já inserido, é desfeita ali. O
procedimento externo recebe 0 e executa o seu \texttt{rollback}, mas não
há mais transação aberta para desfazer. No caminho de sucesso, o exemplo
funciona como esperado. Uma forma de evitar o problema é o procedimento
interno só abrir transação quando \texttt{@@trancount} for 0.

\subsection*{Complemento: procedimentos no padrão SQL}

\textcite[seção 13.4]{Elmasri2011} apresenta os procedimentos
armazenados do padrão SQL/PSM (\emph{Persistent Stored Modules}), com
parâmetros \texttt{in}, \texttt{out} e \texttt{inout} e comandos de
controle de fluxo. Além do reuso por várias aplicações, que a aula
também cita, o livro aponta duas vantagens: menos tráfego entre cliente
e servidor, porque o processamento fica no servidor, e a possibilidade
de verificar regras complexas demais para as restrições declarativas. O
livro observa ainda que muitos SGBDs comerciais aceitam procedimentos
escritos em linguagens de uso geral. O \texttt{output} do T-SQL
corresponde ao \texttt{out} e ao \texttt{inout} do padrão.
